\section{Nexa Drive: Task and System}
\label{sec:system}

\subsection{Functions and actions}
\label{sec:functions}

The assistant controls 24 functions of a simulated car (Table~\ref{tab:functions}; full signatures in
Appendix~\ref{app:functions}). Nineteen correspond to properties of the public Android Automotive
\code{VehiclePropertyIds} specification \citep{android2026vhal} (for example \code{set\_ac} $\leftrightarrow$
\code{HVAC\_POWER\_ON}, \code{HVAC\_TEMPERATURE\_SET}; \code{control\_window} $\leftrightarrow$ \code{WINDOW\_MOVE}), several
properties being merged into one function with a position argument, as in \citet{jafarzadeh2026fixedkeys}. Navigation, nearby
search and media follow Google's documentation for Android car apps, and weather follows CAR-bench. Value ranges
(temperature 16--30\,°C, fan 0--7, volume 0--40) describe one simulated configuration, not a standard. A 25th token,
\code{irrelevant\_function}, means ``refuse''.

\begin{table}[t]
\centering\small
\begin{tabular}{@{}lp{4.9cm}@{}}
\toprule
Group & Functions \\
\midrule
Status (read-only) & fuel level, battery range, tire pressure, vehicle status, odometer, range estimate, weather \\
Climate & AC (on/off, temperature, delta, mode, zone), fan speed, defrost \\
Access & lock, unlock$^{\ast}$, open trunk$^{\ast}$, close trunk, open charge port$^{\ast}$, open fuel flap$^{\ast}$ \\
Windows & side windows, sunroof \\
Lights & headlights, hazard lights \\
Media, route & volume, play media, navigate, find nearby \\
\bottomrule
\end{tabular}
\caption{The 24 functions. $^{\ast}$Sensitive functions (they open access to the car): never proposed from an implicit
request and always confirmed when triggered by the wake word.}
\label{tab:functions}
\end{table}

For every utterance the model chooses one of three actions:
\begin{itemize}
  \item \textbf{Execute} --- \vi{Hạ kính ghế phụ một nửa}{lower the passenger window halfway} $\rightarrow$\\
  {\small\tok{nexa\_15}\code{(window="front\_right",}\\ \code{position="down", amount=50)}\tok{nexa\_end}}.
  \item \textbf{Clarify} --- \vi{Ôi tôi nóng quá}{oh, I am so hot} $\rightarrow$\\
  {\small\tok{nexa\_ask}\tok{nexa\_6}\code{(state="on")}\tok{nexa\_end}}; the runtime asks ``Shall I turn on the air
  conditioning?'' and executes the stored call if the driver says yes, without calling the model again.
  \item \textbf{Refuse} --- \vi{Đặt giúp tôi vé xem phim}{book me a movie ticket} $\rightarrow$
  {\small\tok{nexa\_24}\code{()}\tok{nexa\_end}}.
\end{itemize}
Independent requests in one call are separated by \code{;}. Data policy fixed early in the project: read-only functions are
never ``clarified'' (reading a value cannot do harm, asking only annoys the driver); a clarification keeps only arguments the
user actually stated; sensitive functions and the high beam are never proposed from an implicit request.

\subsection{Functional tokens and prompts}
\label{sec:tokens}

We add 27 tokens to the Qwen2.5 tokenizer \citep{qwen2024qwen25}: \tok{nexa\_0}\ldots\tok{nexa\_24}, \tok{nexa\_end} and
\tok{nexa\_ask} (vocabulary 151,665 $\rightarrow$ 151,692); new embeddings are initialized to the mean of existing ones plus
noise. Because the student is an instruction-tuned model, prompts use its native ChatML template. The call prompt is a
one-line system message (``You are an agent that controls a car'', in Vietnamese) and the user utterance; it averages 33--39
tokens. From data round v6 on, the user turn may start with \code{Lượt trước: <previous call>} (``previous turn''), which lets
the model resolve follow-ups such as \vi{đóng đi}{close it}. As in Octopus~v2, we tested appending a natural-language
description of the function after \tok{nexa\_end} (\code{Mô tả hàm: \ldots}, ``function description''; \S\ref{sec:rounds}).
The final training set keeps it in only 2,385 of its 33,523 examples (7.1\%, refusals inherited from earlier rounds);
at inference the model stops at \tok{nexa\_end}, so the description costs no latency. A second prompt type (the call, the function result and
the user turn) trains the same model to phrase the answer; examples are in Appendix~\ref{app:prompts}.

\subsection{Runtime pipeline}
\label{sec:pipeline}

\begin{figure*}[t]
\centering
\begin{tikzpicture}[
  box/.style={draw, rounded corners=2pt, align=center, font=\scriptsize, minimum height=0.95cm, inner sep=3pt},
  model/.style={box, fill=blue!10},
  code/.style={box, fill=orange!12},
  io/.style={box, fill=green!10},
  arr/.style={-{Latex[length=1.6mm]}, semithick},
  node distance=0.32cm]
\node[io] (mic) {Microphone\\16\,kHz};
\node[model, right=of mic] (wake) {Wake word\\``Hey Nexa''\\or push-to-talk};
\node[model, right=of wake] (stt) {STT\\PhoWhisper-small\\int8, CPU};
\node[code, right=of stt] (norm) {Normalize\\numbers, case;\\split multi-\\command};
\node[model, right=of norm] (llm) {Student 1.5B\\GGUF q4\_K\_M\\\tok{nexa\_i}(args)};
\node[code, right=of llm] (thr) {Confidence\\threshold\\0.80};
\node[code, below=0.55cm of thr] (val) {Parser +\\validator};
\node[code, left=of val] (safe) {Safety rules\\(vehicle state,\\source)};
\node[io, left=of safe] (car) {Vehicle\\simulator};
\node[code, left=of car] (resp) {Template\\answer};
\node[model, left=of resp] (tts) {TTS\\VieNeu};
\node[io, left=of tts] (ui) {Speaker +\\3D interface};
\draw[arr] (mic) -- (wake); \draw[arr] (wake) -- (stt); \draw[arr] (stt) -- (norm);
\draw[arr] (norm) -- (llm); \draw[arr] (llm) -- (thr); \draw[arr] (thr) -- (val);
\draw[arr] (val) -- (safe); \draw[arr] (safe) -- (car); \draw[arr] (car) -- (resp);
\draw[arr] (resp) -- (tts); \draw[arr] (tts) -- (ui);
\end{tikzpicture}
\caption{Runtime pipeline. Blue: learned models; orange: code; green: input/output. Clarifications and refusals leave the
pipeline after the validator or safety step with a question or an apology; a pending call is executed when the driver
confirms. Everything runs on the laptop CPU; the GPU is not used.}
\label{fig:pipeline}
\end{figure*}

Figure~\ref{fig:pipeline} shows the deployed pipeline. The backend (FastAPI + WebSocket) keeps one simulated vehicle and one
pending call per session. The student runs in \code{llama-server} from llama.cpp \citep{gerganov2023llamacpp}; prompts are
sent as token ids produced by the Hugging Face tokenizer and outputs are decoded from token ids, because the text field of
the server hides control tokens such as \tok{nexa\_i}. The end-of-sequence token is overridden to \tok{nexa\_end}: without it,
uncertain refusals continued into the description seen in those 7.1\% of examples until the generation limit, which made them five to eight times
slower (identical outputs on 400 validation prompts, total time 540\,s $\rightarrow$ 249\,s).

Several small layers sit around the model; each was added after a measured failure (details in Appendix~\ref{app:runtime}):
\begin{itemize}
  \item \textbf{Text normalization.} The training data, written by the teacher, always starts with a capital letter. For
  \vi{hỏng xe giữa đường rồi}{my car broke down in the middle of the road}, the lower-case version is refused (refusal
  probability 0.68) while the capitalized one triggers a hazard-light proposal (0.99). We therefore capitalize the first
  letter, collapse repeated letters and lower-case all-caps input before the model sees it.
  \item \textbf{Retry without context.} If the model refuses while a previous turn is in the prompt, it is called once more
  without the context, so that context can help a follow-up but never turns an understood request into a refusal.
  \item \textbf{Relative-command conventions.} ``Colder'' without a number means $-2$\,°C, fan $\pm1$ level, volume $\pm3$---the
  same rules used to label the training data.
  \item \textbf{Multi-command splitting.} The training data contain no example that executes several independent actions, so
  the model picks one action from \vi{bật đèn pha lên xong kiểm tra xăng xong đó mở cốp}{turn on the headlights, then check
  the fuel, then open the trunk}. We split such utterances on explicit sequential connectors (\textit{xong}, \textit{sau
  đó}, \textit{rồi thì}, \ldots; not on \textit{và} ``and''), run the model per segment and stop at the first segment that
  needs clarification.
  \item \textbf{Answer templates.} The model was trained to phrase answers from function results, but this second pass adds
  about 1.6\,s and produced semantic errors no automatic check caught (\S\ref{sec:turn2}); the deployed system answers with
  templates filled from the result.
\end{itemize}

\subsection{Guardrails}
\label{sec:guardrails}

Following the principle ``the model understands, code keeps the car safe'', every call passes a validator and a set of
state-dependent rules declared in the function schema (Table~\ref{tab:safety}). The validator rejects values outside
enumerations or ranges, asks for missing required arguments and drops unknown arguments. Each command carries its source:
push-to-talk is treated as the driver, while the wake word can be spoken by any passenger. Mirroring a real window-lock
button, the child lock lets the driver lower rear windows directly but asks for confirmation when the request comes through
the wake word. The rules are covered by 917 automated tests, including scenario tests that, as in $\tau$-bench, check the
final vehicle state rather than the wording of the answer. We did not measure guardrails with a quantitative benchmark.

\begin{table}[t]
\centering\small
\begin{tabular}{@{}p{3.6cm}p{3.1cm}@{}}
\toprule
Condition & Action \\
\midrule
open trunk / charge port / fuel flap while driving & refuse \\
unlock doors while driving & ask to confirm \\
any sensitive function via wake word & ask to confirm \\
child lock on + wake word + rear window down or sunroof open/tilt & ask to confirm \\
driving starts with trunk, door or flap open & warn \\
model proposes a sensitive function & drop the proposal \\
\bottomrule
\end{tabular}
\caption{Runtime safety rules (code, independent of the model).}
\label{tab:safety}
\end{table}
